\section{Discussion}

\benchmarkname{} is best understood as a preflight benchmark for public-facing
AI systems. It does not ask whether a model can solve the hardest available
questions. It asks whether the model can avoid small mistakes that users
reasonably expect it never to make.

The 2026-06-03 evidence snapshot has two useful readings. First, a few frontier
and well-routed settings now pass this small public split completely: three
model settings answer all 224 questions correctly. Second, failures remain highly
legible when they occur. The benchmark surfaces counting mistakes, spelling
transform mistakes, constraint slips, verbose answers to short-answer prompts,
provider refusals, and cost/token tradeoffs in a form that product teams can
inspect directly.

The narrow scope is the point. Short objective prompts make deterministic
scoring, cheap repeated runs, and failure galleries practical. That makes the
benchmark useful for model comparison, prompt regression testing, and product
quality monitoring. The right use is not to crown a single model from a small
public split, but to track whether obvious failures recur across versions,
settings, and providers.

The practical recommendation is to track obvious failure modes explicitly:
separate answer correctness from format compliance, preserve examples that
product teams can understand, report provider errors as counted attempts, and
refresh the set over time rather than optimizing to one static leaderboard.
